\section{Gaussian domination on coordinate tori}\label{sec:coordinate}

The compactness estimate comes from the coordinate trace. We first
retain the real slice connections and all magnetic fluxes in that
trace, then remove the external curvature field by two applications
of Schatten H\"older in different time directions. We give the all-flux
trace and Gaussian-domination proof of
\cite[Section~5.8 and Theorem~6.27]{ChamberQED}.

\subsection{Coordinate slices and the all-flux trace}\label{subsec:slab-coordinate-trace}

For $s>0$ write
\[
 g_s(x)=(2\pi s)^{-1/2}e^{-x^2/(2s)}.
\]
Tensor powers of $g_s$ denote product densities. Thus $e^{-a|E|^2}$
has convolution kernel $g_{2a}$ in each real link variable.

For Gaussian domination we use a slice representation with real
spatial connections. Take direction $1$ as time on
$C_{\mathbf n}=C(n_1,n_2,n_3)$, $\mathbf n=(n_1,n_2,n_3)$,
$n_\mu\ge4$, and let
$\Sigma=\Z/n_2\Z\times\Z/n_3\Z$. Its links and plaquettes are
$\Sigma^{(1)}$ and $\Sigma^{(2)}$, its coboundary is $d_s$, and
$v=|\Sigma|=n_2n_3$. A connection has spatial slices $A_j$ and
temporal links $a_j$, and its curvature has the parts
\begin{equation}\label{eq:slab-slice-curvatures}
 F_j^{\mathrm{el}}=A_{j+1}-A_j-d_sa_j+r_qz_j^{\mathrm{el}},
 \qquad F_j^{\mathrm{mag}}=d_sA_j+r_qz_j^{\mathrm{mag}}.
\end{equation}
Closedness reads
$z_{j+1}^{\mathrm{mag}}-z_j^{\mathrm{mag}}=d_sz_j^{\mathrm{el}}$.

We first unfold the integer variables. Put
$k_0=0$, $k_{j+1}=k_j+z_j^{\mathrm{el}}$. Closedness is equivalent to
\[
 z_j^{\mathrm{mag}}=z_0^{\mathrm{mag}}+d_sk_j,
                       \qquad d_sk_{n_1}=0.
\]
The integer shift with spatial values $k_j$ and zero temporal values
removes all electric integers except the last one, $k_{n_1}$, and
changes every spatial slice to $A_j+r_qk_j$. Summing $k_j$ while
integrating $A_j$ over its compact box unfolds the intermediate
slices to all of $\R^{\Sigma^{(1)}}$. The last integer remains as a
winding vector
\[
 g=r_qk_{n_1}\in\Gamma_s,
 \qquad\Gamma_s=r_q(\Z^{\Sigma^{(1)}}\cap\ker d_s).
\]

Fix a plaquette $\sigma_0$ of $\Sigma$ and write
$k=\sum_\sigma z_0^{\mathrm{mag}}(\sigma)$. Then
$z_0^{\mathrm{mag}}-k\mathbf1_{\sigma_0}=d_sk'$ for an integer
one-cochain $k'$. Integer exactness follows from the dual graph whose
vertices are the plaquettes. A link of
$\Sigma$ is a dual edge directed toward the plaquette where its
coefficient in $d_s$ is positive. The dual graph is connected. Fix
a root and a spanning tree. Unit signed flow along the path from the
root to a vertex $\sigma$ has divergence
$\delta_\sigma-\delta_{\sigma_0}$. Summing these paths with the
integer coefficients of a zero-sum two-cochain gives precisely that
cochain as its divergence, hence as $d_sk'$. Two choices of $k'$
differ by an integer element of $\ker d_s$.

Shift every slice by $r_qk'$. The magnetic integers now equal
$k\mathbf1_{\sigma_0}$ in every slice. Summation over $k'$ modulo
closed integer cochains replaces the compact box for $A_0$ by a
fundamental domain of $\Gamma_s$. We arrive at paths
\[
 A_0\in\R^{\Sigma^{(1)}}/\Gamma_s,\qquad
 A_j\in\R^{\Sigma^{(1)}}\ (0<j<n_1),\qquad
 A_{n_1}=A_0+g,\quad g\in\Gamma_s,
\]
with $a_j\in[0,r_q)^{\Sigma}$ and flux $k\in\Z$. Their curvatures are
\begin{equation}\label{eq:slab-unfolded-curvatures}
 F_j^{\mathrm{el}}=A_{j+1}-A_j-d_sa_j,
 \qquad F_j^{\mathrm{mag}}=d_sA_j+r_qk\mathbf1_{\sigma_0}.
\end{equation}
Thus integration over the original configuration quotient equals
summation over $k,g$ and integration over these paths. For a
nonnegative integrand the identity follows by Tonelli's theorem. For
an absolutely integrable complex integrand it follows by applying
that identity to its modulus and then using Fubini's theorem. Each
change above preserves link phases and curvature, so it applies to
the determinant and to all curvature observables used here.

There are two finite Fock spaces in this representation: $\mathcal F'$
at a coordinate site plane and $\mathcal F$ at a link plane. The
matrix of the open half-slab gives a map
$N(A):\mathcal F'\to\mathcal F$, and we include both crossing weights
by writing
\[
 \widehat N(A)=\omega^{1/2}N(A)(\omega')^{1/2}.
\]
Let $u(a)$ be the diagonal unitary on $\mathcal F$ with entries
$e^{iq\langle a,\mathfrak q(I)\rangle}$. The fermion trace is
\begin{equation}\label{eq:slab-slice-determinant}
 \det \mathcal D_A^{C_{\mathbf n}}=
    \tr_{\mathcal F'}\prod_{j=n_1-1,\ldots,0}
             \bigl(\widehat N(A_{j+1})^*u(a_j)\widehat N(A_j)\bigr).
\end{equation}
This formula also holds on an unfolded path, since
$A_{n_1}-A_0\in\Gamma_s$ leaves $\widehat N$ unchanged.

To derive \eqref{eq:slab-slice-determinant}, split the torus at each
site and link plane. The half-slab above a link plane is the reflected
half-slab below it, hence its matrix is $N(A_{j+1})^*$. The crossing
at a site plane has no link phase and gives $\omega'$. At a link
plane it gives $\omega^{1/2}u(a_j)\omega^{1/2}$. Compose the matrices
and close antiperiodically; the closing parity is the identity.
Cyclicity moves half of the final $\omega'$ to the front, giving
\eqref{eq:slab-slice-determinant}. Every entry of $\widehat N(A)$ is
a polynomial in the link phases: the half-slab action is quadratic
in finitely many generators, its exponential terminates, and
Berezin integration extracts a coefficient. Consequently
\begin{equation}\label{eq:slab-half-matrix-bound}
 C_N:=\sup_{A\in\R^{\Sigma^{(1)}}}\|\widehat N(A)\|<\infty,
 \qquad\widehat N(A+r_qk')=\widehat N(A)
       \quad(k'\in\Z^{\Sigma^{(1)}}).
\end{equation}

On real slice connections the two Hilbert spaces are
\[
 H^s=\bigoplus_{k\in\Z}
      L^2(\R^{\Sigma^{(1)}}/\Gamma_s;\mathcal F'),\qquad
 H^l=\bigoplus_{k\in\Z}
      L^2(\R^{\Sigma^{(1)}}/\Gamma_s;\mathcal F).
\]
All quotient measures are induced by Lebesgue measure. Define commuting
self-adjoint electric operators $E_b$ by
$e^{-i\langle x,E\rangle}\Phi(A)=\Phi(A-x)$. On $H^l$ let
\begin{equation}\label{eq:slab-electric-operators}
 U_G(a)\Phi(A)=u(a)\Phi(A-d_sa),\qquad
 P_G=r_q^{-v}\int_{[0,r_q)^{\Sigma}}U_G(a)\,\dd a,
 \qquad K=r_q^ve^{-|E|^2/2}P_G.
\end{equation}
The integer mode charges make $U_G$ a representation of the compact
gauge group $\R^{\Sigma}/r_q\Z^{\Sigma}$. Its average $P_G$ is the
orthogonal projection onto invariant vectors. It commutes with every
function of $E$, because $u(a)$ acts on the Fock index and all spatial
translations commute. Hence
\begin{equation}\label{eq:slab-electric-square-root}
 K^{1/2}=r_q^{v/2}e^{-|E|^2/4}P_G.
\end{equation}

For a real slice two-cochain $b$, define the multiplication operator
$M_b:H^s\to H^l$ by
\begin{equation}\label{eq:slab-magnetic-multiplier}
 (M_b\Phi)_k(A)=
     e^{-\|d_sA+r_qk\mathbf1_{\sigma_0}+b\|^2/4}
                                      \widehat N(A)\Phi_k(A),
 \qquad X_b=K^{1/2}M_b.
\end{equation}
The magnetic weight is bounded and $\widehat N$ is bounded, so $M_b$
is bounded. We must prove that $X_b$ is Hilbert--Schmidt; the real
quotient is not compact and there are infinitely many flux sectors.

The periodized Gaussian
\[
 \widetilde g_{1/2}([x])
   =\sum_{\gamma\in\Gamma_s}\prod_{b\in\Sigma^{(1)}}
                       \pi^{-1/2}e^{-(x_b+\gamma_b)^2}
\]
has integral one on the quotient and is bounded. Indeed
$\Gamma_s\subset r_q\Z^{\Sigma^{(1)}}$, so the series is dominated
by the product of one-dimensional Gaussian sums on that larger lattice.
It is therefore in $L^1\cap L^\infty\subset L^2$. The kernel of
$K^{1/2}$ is
\[
 K^{1/2}(A',A)=r_q^{-v/2}\int_{[0,r_q)^{\Sigma}}
       \widetilde g_{1/2}([A'-A-d_sa])u(a)\,\dd a.
\]
Its operator norm is bounded by the nonnegative function
\[
 \widehat g([x])=r_q^{-v/2}\int_{[0,r_q)^{\Sigma}}
                  \widetilde g_{1/2}([x-d_sa])\,\dd a,
\]
which has finite $L^1$ and $L^\infty$ norms and hence finite $L^2$ norm.
Integrating the square of the kernel of $X_b$ over $A'$ now gives
\begin{equation}\label{eq:slab-shifted-hs}
 \|X_b\|_2^2\le
    (\dim\mathcal F')C_N^2\|\widehat g\|_2^2
       \sum_{k\in\Z}\int_{\R^{\Sigma^{(1)}}/\Gamma_s}
             e^{-\|d_sA+r_qk\mathbf1_{\sigma_0}+b\|^2/2}\,\dd A.
\end{equation}

To bound this last sum, split the quotient as
\[
 (\ker d_s/\Gamma_s)\times(\ker d_s)^\perp.
\]
The first factor is compact: $\ker d_s$ is the kernel of an integer
matrix, hence has an integer basis, and $\Gamma_s$ has full rank in
it. On the second factor $d_s$ is injective, so
$\|d_sA\|^2$ is positive definite. Moreover
$\ker d_s^*$ consists of constant plaquette functions. This follows
from the two equations
\[
 (d_s^*b)(y,2)=b(y)-b(y-e_3),\qquad
 (d_s^*b)(y,3)=b(y-e_2)-b(y).
\]
Thus $\operatorname{ran}d_s$ is the orthogonal complement of the
constant two-cochain. The orthogonal constant component of
$r_qk\mathbf1_{\sigma_0}+b$ has squared norm
$v^{-1}(r_qk+\sum_\sigma b_\sigma)^2$. Translate the variable in
$(\ker d_s)^\perp$ to remove the remaining component. We obtain
\begin{equation}\label{eq:slab-flux-gaussian}
 \sum_k\int e^{-\|d_sA+r_qk\mathbf1_{\sigma_0}+b\|^2/2}\,\dd A
    =C_s\sum_{k\in\Z}
               e^{-(r_qk+\sum_\sigma b_\sigma)^2/(2v)}<\infty,
\end{equation}
where
$C_s=\operatorname{vol}(\ker d_s/\Gamma_s)
\int_{(\ker d_s)^\perp}e^{-\|d_sA\|^2/2}\,\dd A<\infty$.
This constant is independent of $k$ and $b$. Equations
\eqref{eq:slab-shifted-hs}--\eqref{eq:slab-flux-gaussian} prove the
Hilbert--Schmidt property for every real magnetic perturbation. The
bounds here depend on the fixed cross-section; their purpose is to
justify the trace argument.

For a real external two-cochain $\beta$ on $C_{\mathbf n}$, put
\begin{equation}\label{eq:slab-external-partition}
 Z_{C_{\mathbf n}}(\beta)=\int_{\operatorname{Conf}(C_{\mathbf n})}
          e^{-\|F+\beta\|^2/2}\det \mathcal D_A^{C_{\mathbf n}}\,\dd(A,z).
\end{equation}
Let $\beta_j^{\mathrm{el}}$ and $\beta_j^{\mathrm{mag}}$ be its
electric and magnetic slices. Then
\begin{equation}\label{eq:slab-external-trace}
 Z_{C_{\mathbf n}}(\beta)=(2\pi)^{n_1n_2n_3}
       \operatorname{Tr}_{H^s}\prod_{j=n_1-1,\ldots,0}
            (X_{j+1}^*U_jX_j),
 \quad X_j=X_{\beta_j^{\mathrm{mag}}},\quad
 U_j=e^{i\langle\beta_j^{\mathrm{el}},E\rangle},
 \quad X_{n_1}=X_0.
\end{equation}

Here is the conversion from the unfolded path integral. The electric
Maxwell factors are the kernels
\[
 e^{-|A_{j+1}-A_j-d_sa_j+\beta_j^{\mathrm{el}}|^2/2}
       =(2\pi)^v g_1^{\otimes 2v}
                       (A_{j+1}-A_j-d_sa_j+\beta_j^{\mathrm{el}}).
\]
After summing the winding vector they are periodized on $\Gamma_s$.
Gaussian convolution with a translated kernel is
$e^{i\langle\beta_j^{\mathrm{el}},E\rangle}e^{-|E|^2/2}$.
Integration over $a_j$ includes the matrix $u(a_j)$ from
\eqref{eq:slab-slice-determinant}; it is exactly the gauge average in
\eqref{eq:slab-electric-operators}. The magnetic Gaussian in slice
$j$ splits into the two quarter-Gaussians in $M_j$ and $M_j^*$.
The determinant half-slab matrices supply their other factors.
Thus the cyclic path integral is the trace on $H^l$ of
\[
 \prod_{j=n_1-1,\ldots,0}
       (K^{1/2}U_jK^{1/2}M_jM_j^*).
\]
Each time slice contributed $(2\pi)^v$, giving the constant in
\eqref{eq:slab-external-trace}. Move the final $M_0^*$ to the front
and regroup the products as $X_{j+1}^*U_jX_j$. The
Hilbert--Schmidt property above and the absolute Gaussian path
majorant justify this cyclicity and every interchange. In particular
the external-field integral is finite; alternatively this follows
directly by absorbing a linear factor into a smaller Gaussian.

We spell out the path and kernel identification used here. On a
fundamental domain write every intermediate real slice as
$A_j=A_j^\flat+g_j$, $g_j\in\Gamma_s$, and set $g_0=0$ and
$g_{n_1}=A_{n_1}-A_0$. Each unperiodized electric kernel then has
the argument
$A_{j+1}^\flat-A_j^\flat+g_{j+1}-g_j-d_sa_j$.
The map $(g_1,\ldots,g_{n_1})\mapsto
(g_{j+1}-g_j)_{0\le j<n_1}$ is a bijection of $\Gamma_s^{n_1}$.
Summing these differences independently is therefore precisely
periodization of every cyclic kernel. The magnetic factors and
$\widehat N$ are unchanged by the shifts $g_j$. The finite Fock trace
is changed from $\mathcal F'$ to $\mathcal F$ by moving its last
$\widehat N(A_0)^*$ to the front; periodicity identifies it with
$\widehat N(A_{n_1})^*$.

For absolute convergence replace every matrix factor by its norm,
and bound each shifted electric Gaussian by a fixed constant times
$e^{-|A'-A|^2/4}$ and each shifted magnetic factor by a constant
times $e^{-\|F^{\mathrm{mag}}\|^2/4}$. This defines nonnegative
scalar kernels. Their periodized convolutions are bounded and in
$L^2$, and multiplication by the magnetic Gaussian is square
integrable after summing fluxes by \eqref{eq:slab-flux-gaussian}.
Hence each scalar majorant operator is Hilbert--Schmidt. The
integral of their cyclic product is finite: group the first factor
against the product of the others and use the Hilbert--Schmidt
Cauchy--Schwarz inequality. Tonelli first gives this conclusion
for the nonnegative majorants, and Fubini then permits the kernel
expansions and periodization changes for the original integrand.
This proves the trace identity without a formal interchange of
the infinite flux sum and the operator trace.

\subsection{Gaussian domination}\label{subsec:slab-gaussian}

\begin{proposition}\label{thm:slab-gaussian}
For every coordinate torus with sides at least four and every real
two-cochain $\beta$,
\begin{equation}\label{eq:slab-gaussian-partition}
 0<Z_{C_{\mathbf n}}(\beta)\le Z_{C_{\mathbf n}}(0).
\end{equation}
Consequently its normalized curvature law obeys
\begin{equation}\label{eq:slab-gaussian-domination}
 \E_{C_{\mathbf n}}e^{\langle\beta,F\rangle}
                              \le e^{\|\beta\|^2/2}.
\end{equation}
For real two-cochains $\beta_i$ and complex numbers $z_i$,
\begin{equation}\label{eq:slab-complex-gaussian}
 \left|\E_{C_{\mathbf n}}e^{\sum_i z_i\langle\beta_i,F\rangle}\right|
       \le\exp\left(\tfrac12
                    \left\|\sum_i(\operatorname{Re}z_i)\beta_i\right\|^2\right).
\end{equation}
The constants are independent of the volume, the charge, and the mass.
\end{proposition}

\begin{proof}
The determinant is nonnegative and positive on a set of positive
configuration measure, since $Z_{C_{\mathbf n}}(0)>0$. The shifted photon
density is strictly positive everywhere, proving the strict lower
bound in \eqref{eq:slab-gaussian-partition}.

In \eqref{eq:slab-external-trace} move $X_0^*$ through the trace to
obtain a trace on $H^l$:
\[
 (2\pi)^{-n_1n_2n_3}Z_{C_{\mathbf n}}(\beta)
       =\operatorname{Tr}_{H^l}
          (U_{n_1-1}X_{n_1-1}X_{n_1-1}^*\cdots U_0X_0X_0^*).
\]
Every $U_j$ is unitary. Apply Schatten's inequality with all exponents
equal to $n_1$; the trace on the left is nonnegative. Since $X_jX_j^*$
and $X_j^*X_j$ have the same nonzero eigenvalues, this gives
\begin{equation}\label{eq:slab-magnetic-reduction}
 Z_{C_{\mathbf n}}(\beta)
       \le\prod_{j=0}^{n_1-1}
            \bigl((2\pi)^{n_1n_2n_3}
                    \operatorname{Tr}(X_j^*X_j)^{n_1}\bigr)^{1/n_1}
       =\prod_{j=0}^{n_1-1}Z_{C_{\mathbf n}}(\beta^{(j)})^{1/n_1}.
\end{equation}
Here $\beta^{(j)}$ has zero electric part and the same magnetic
part $\beta_j^{\mathrm{mag}}$ in every time slice. The equality uses
the trace formula again, with the same $X_j$ in each slice and all
unitary factors equal to the identity.

Each $\beta^{(j)}$ is supported only on $23$-plaquettes. Interchange
directions $1$ and $2$. The torus becomes $C(n_2,n_1,n_3)$ and the
perturbation is supported only on $13$-plaquettes, so its magnetic
part for the new time direction is zero. The reduction
\eqref{eq:slab-magnetic-reduction} on this torus bounds its partition
function by the unperturbed one. Lattice isometry covariance of the
determinant and photon measure identifies these partition functions
with $Z_{C_{\mathbf n}}(\beta^{(j)})$ and $Z_{C_{\mathbf n}}(0)$. Substitution in
\eqref{eq:slab-magnetic-reduction} proves the upper bound in
\eqref{eq:slab-gaussian-partition}.

Replacing $\beta$ by $-\beta$ and completing the square gives
\eqref{eq:slab-gaussian-domination}. For complex parameters, take
the absolute value inside the probability integral and apply that
bound to $\sum_i(\operatorname{Re}z_i)\beta_i$, proving
\eqref{eq:slab-complex-gaussian}.
\end{proof}


\subsection{Coordinate limits and the reflection form}\label{subsec:coordinate-limits}

\begin{proof}[The coordinate part of \Cref{thm:main}]
The determinant and normalization results make $\mu_n$ a probability
law on $\Omega$ for every integer $n\ge4$. If a finitely supported
cochain $\beta$ embeds in the torus without overlap, its
periodization $\beta_n$ has
\[
 \langle\beta_n,F\rangle=\langle\beta,\widetilde F\rangle,
 \qquad \|\beta_n\|=\|\beta\|.
\]
\Cref{thm:slab-gaussian} therefore gives the claimed Gaussian
bound for this test under $\mu_n$. Every fixed test embeds once
$n$ is large. A single plaquette embeds for every $n\ge4$: its
box of radius $1/2$, enlarged by $[-1,1]^3$, has side three.
Applying the bound with $\beta=\pm R1_\sigma$ and then Markov's
inequality gives
\[
 \mu_n\{|F_\sigma|>R\}\le2e^{-R^2/2},\qquad R>0,
\]
uniformly in $n\ge4$.

Enumerate the plaquettes as $\sigma_1,\sigma_2,\ldots$. Given
$0<\varepsilon<1$, put
\[
 R_j=\sqrt{2\log(2^{j+1}/\varepsilon)},\qquad
 K_\varepsilon=\Omega\cap\prod_{j\ge1}[-R_j,R_j].
\]
The set $\Omega$ is closed in the countable product, so
$K_\varepsilon$ is compact. The union bound gives
$\mu_n(\Omega\setminus K_\varepsilon)\le\varepsilon$ for all
$n\ge4$. Thus the family is tight. Since $\Omega$ is Polish,
Prokhorov's theorem gives subsequential weak limits; see
\cite[Theorem~1.1]{Sorensen}.

Let $\mu_{n_i}\to\nu$ along such a subsequence, and fix $\beta$.
For all sufficiently large $i$, the support of $\beta$ embeds.
For $K>0$ the function
$\min\{e^{\langle\beta,F\rangle},K\}$ is bounded and continuous.
Weak convergence followed by monotone convergence in $K$ gives
\textup{(i)}. Each $\mu_{n_i}$ is invariant under $W$ and charge
conjugation by the finite-volume symmetry identities. The cubic
period lattice and its character, equal to $-1$ on each coordinate
period, are preserved by every signed coordinate permutation.
Translations preserve them as well. Their pullbacks and $F\mapsto-F$
are continuous on $\Omega$, so these symmetries pass to $\nu$.
This proves \textup{(ii)} and reflection invariance at every mirror.

Fix a coordinate mirror $\Pi=\{x_j=c\}$, $c\in\tfrac12\Z$,
and a bounded continuous function $G$ of finitely many plaquettes
wholly in $\{x_j>c\}$. Choose $h_G>0$ so that every corner of
these plaquettes satisfies $c<x_j<c+h_G$. For $n_i>2h_G$, the
support lies in the open positive half-period slab
\[
 \{c<x_j<c+n_i/2\}.
\]
The cube is adapted to this mirror, with closing vector $n_ie_j$
and character $-1$. Its other seam is at $c+n_i/2$, which belongs
to $\tfrac12\Z$ for both even and odd $n_i$. Thus both seams are
admissible, including when one is a site plane and the other a
link plane. The adapted positivity theorem applies and gives
\[
 \mu_{n_i}\bigl(\overline{G(\theta_\Pi^*F)}G(F)\bigr)\ge0.
\]
The integrand is bounded and continuous, so the inequality passes
to $\nu$.

\label{par:side-density}
We give the extension to the whole side, which will also be used
for the diagonal mirrors. Suppose a law $\lambda$ is invariant
under $\theta_\Pi$ and its reflected form is nonnegative on
bounded continuous local tests on one open side. On each
finite-coordinate marginal, compactly supported continuous
functions are dense in $L^2$. Reflection invariance and
Cauchy--Schwarz give
\[
 \left|\lambda\bigl((\Theta A)A-(\Theta B)B\bigr)\right|
 \le\|A-B\|_{L^2(\lambda)}
       \bigl(\|A\|_{L^2(\lambda)}+\|B\|_{L^2(\lambda)}\bigr),
 \qquad \Theta A=\overline{A\circ\theta_\Pi^*}.
\]
Hence positivity extends to all finite-coordinate $L^2$ tests.
Enumerate the plaquettes in the chosen open side, and let
$\mathcal F_m$ be generated by the first $m$. For
$A\in L^2(\Sigma_H,\lambda)$, the conditional expectations
$\lambda(A\mid\mathcal F_m)$ converge to $A$ in $L^2$, because
$\Sigma_H$ is generated by $\bigcup_m\mathcal F_m$. The same
inequality extends positivity to $A$. Reflection carries one
open side to the other, giving the assertion on either side.
Applied to $\lambda=\nu$, this proves the coordinate part of
\textup{(iii)}. The diagonal part remains to be proved by comparison.
\end{proof}

\medskip\noindent
\emph{Optimality of the Gaussian constant.}
Choose a nonzero finitely supported integer two-cochain $z_0$ with
$dz_0=0$; for example, take $z_0=dk$ with $k$ equal to one on a
single oriented link and zero on the other links. Put $\beta=r_qz_0$.
For all sufficiently large $n$, these cochains embed in the cubic
torus. We use the same symbols for their periodizations, which
preserve $dz_0=0$, the source norm, and the pairing with the
periodically extended field. Completing the square gives
\[
\begin{split}
 &\sum_{dz=0}\int
 e^{-\|dA+r_qz\|^2/2+\langle\beta,dA+r_qz\rangle}
       \det\mathcal D_A^{C(n)}\,\dd A\\
 &\qquad =
 e^{\|\beta\|^2/2}\sum_{dz=0}\int
 e^{-\|dA+r_q(z-z_0)\|^2/2}
       \det\mathcal D_A^{C(n)}\,\dd A
 =e^{\|\beta\|^2/2}Z_{C(n)}[1].
\end{split}
\]
Here the sum runs over all closed integer two-cochains and the
integrals have the same domain as in the definition of $\mu_n$.
The last equality reindexes this entire sum; the determinant
depends only on $A$. Thus
$\mu_n(e^{\langle\beta,F\rangle})=e^{\|\beta\|^2/2}$.
Put $X(F)=e^{\langle\beta,F\rangle}$. The same identity for
$2\beta$ gives $\mu_n(X^2)=e^{2\|\beta\|^2}$ for large $n$,
and \textup{(i)} gives $\nu(X^2)\le e^{2\|\beta\|^2}$.
For either of these laws $\lambda$,
\[
 0\le\lambda(X)-\lambda(X\wedge K)
 \le K^{-1}\lambda(X^2)\le K^{-1}e^{2\|\beta\|^2}.
\]
The function $X\wedge K$ is bounded and continuous. First passing
to the weak limit at fixed $K$, then letting $K\to\infty$, gives
\[
 \nu(e^{\langle\beta,F\rangle})=e^{\|\beta\|^2/2}
 \qquad(\beta=r_qz_0).
\]
Since $\beta\ne0$, no constant smaller than one can replace the
constant in \textup{(i)} at the fixed parameters of the theorem.
